\begin{enumerate}[start=7]
\item  L'équation différentielle $y' = -0,1y$ peut s'écrire $y' + 0,1y = 0$.
    
Il s'agit d'une équation différentielle linéaire du premier ordre de la forme :
\[y' + ay = b \text{ avec } a = 0,1 \text{ et } b = 0.\]

La solution générale est :
\[A(t) = k\e^{-at} = k\e^{-0,1t} \text{ où } k \in \mathbb{R}.\]

Avec la condition initiale :
\begin{align*}
A(0) &= \np{7400} \\
k\e^{-0,1\times0} &= \np{7400} \\
k\e^0 &= \np{7400} \\
k &= \np{7400}.
\end{align*}

La solution est donc bien :
\[A(t) = \np{7400}\e^{-0,1t}.\]

\item Au bout de 70 jours, l'activité radioactive vaut environ :
\[A(70)=\np{7400}\e^{-0,1\times 70}=\np{7400}\e^{-7}\approx \np{6,75}.\]

Il reste donc environ $\np{6,75}$ MBq sur les $\np{7400}$ MBq initiaux.

Perdre plus de $99,9\,\%$ de sa valeur initiale signifie qu'il reste moins de $0,1\,\%$ de cette valeur, soit : $0,001\times \np{7400}=\np{7,4}$.

Comme $\np{6,75}<\np{7,4}$, l'activité a bien diminué de plus de $99,9\,\%$ au bout de 70 jours.

\item D'après la question précédente, au bout de 70 jours, l'activité radioactive du lutécium 177 est inférieure à $0,1\,\%$ de sa valeur initiale. Elle est donc considérée comme négligeable et sans risque pour l'environnement.

La préconisation du fabricant est donc cohérente : attendre 70 jours avant de jeter le sac aux ordures ménagères permet de laisser l'activité radioactive diminuer suffisamment.
\end{enumerate}

\bigskip
